\section{Our Approach}
\label{sec:approach}

% Owner: @Capheuss (parent). Only the owner edits this file; send input to them on #40. Shared framing only; per-track text is in 4a-4d.
% Each paragraph below maps to the GitHub issues noted in its comment;
% the overall roadmap is issue #17 and project-docs/PLAN.md.

We implement Speculative Taint Tracking (STT)~\cite{Yu:STT:2019}, followed by
the two $\mu$STT mechanisms~\cite{Chen:uSTT:2025}, in the RTL of SonicBOOM
(BOOMv3)~\cite{Zhao:SonicBOOM:2020}, the out-of-order RV64GC core of the
Chipyard SoC framework~\cite{Amid:Chipyard:2020}. Prior to this work, $\mu$STT had
been evaluated only in gem5 on ARM, and its released RTL is a standalone
rename-stage circuit rather than part of a core. Our contribution is to build
both mechanisms into a complete, synthesizable RISC-V core and measure their
cost. The work was divided into 4 areas: taint tracking, delaying
transmitters at issue, security and performance evaluation, and
hardware cost.

\subsection{STT in BOOM terms}
STT treats a load that executes speculatively as an \emph{access}
instruction and sets the register it writes as \emph{tainted}. The taint then spreads
to every instruction that reads that register or any other tainted register.
A \emph{transmitter} is an instruction whose timing can reveal its operands.
In BOOM these are loads, stores and atomics (through their address and their data in the 
cases of stores), as well as conditional branches and \texttt{jalr} (through their
outcome). A tainted transmitter may not execute until its taint is cleared,
and a taint is cleared only once the load that caused it passes the
\emph{visibility point} (VP), after which it can no longer be squashed. Under
the Spectre threat model (Section~\ref{sec:threat}), the VP is the oldest
unresolved branch in the reorder buffer (ROB). Each register's taint is stored
as the ROB index of the youngest access load it depends on, its
\emph{youngest root of taint} (YRoT). An instruction is safe once its YRoT is
strictly older than the VP; a YRoT equal to the VP is unsafe.

\subsection{Building the mechanism}

% Issue #3 (A1): interface spec, shared by tracks A and B.
\paragraph{Interface}
The interface between the taint tracking and the delaying mechanisms is a written specification.
Each micro-op gains a YRoT (a valid bit and a ROB index) and a transmitter flag, set at 
decode from the original instruction so that split stores keep it. The ROB exports the VP, 
and a shared helper, \textit{isSafe(YRoT)}, compares a YRoT with the VP in
circular ROB order relative to the head. This insures that the check stays correct when the
ROB wraps. The core also exports a recovery signal, \texttt{stt\_block}, which
holds every transmitter during a mispredict, flush or rollback. The delaying mechanism
combines these into the permission each transmitter must obtain before it issues.

Five switches set the
\texttt{conservativeSTT}, \texttt{enableSTT}, \texttt{enableAgeMatrix},
\texttt{impedeUop}, and \texttt{threatModel} mechanisms at elaboration time. 
All are off by default, and illegal combinations fail elaboration, for example the 
Age Matrix without STT, or the conservative stand-in together with STT. The new 
micro-op fields and ports exist only when a switch needs them, so with every switch 
off the generated Verilog matches the unmodified core (Section~\ref{sec:verification}). 
Every measurement therefore compares a protected core with the same baseline.

% Owner: @Capheuss (track A). Only the owner edits this file; send input to them on #40. Taint propagation and the Age Matrix.

% Issue #4 (A2): baseline STT taint propagation.
\paragraph{Taint propagation}
Taint is kept in two tables, one for integer and one for floating-point
physical registers, each holding one YRoT per register, in the style of BOOM's
busy table. The tables sit at the boundary between rename and dispatch rather
than inside the rename stage, because BOOM assigns each micro-op its ROB index
in the core, after rename. An instruction's YRoT is the youngest valid taint
among all of its register sources, including store data and the third operand
of fused multiply-adds; integer \texttt{x0} and immediates contribute nothing.
A load or atomic then writes its own ROB index to its destination, becoming a
new root, while any other instruction writes its YRoT. A load's own YRoT
covers only its address, so a load with an untainted address may still execute
under an unresolved branch. Instructions dispatched in the same cycle must see
each other's taint, so a bypass passes each lane's destination taint to the
younger lanes that read it. Taint tables are written only for lanes that
actually dispatch, so a stalled lane creates no root and recomputes its taint
when it dispatches later.

\paragraph{Visibility point and recovery}
The ROB keeps one pending bit per entry, set when a conditional branch or
\texttt{jalr} enters the ROB and cleared when the branch resolves or the entry
is killed, flushed or committed. The VP is the oldest entry whose bit is set,
found with BOOM's existing age-priority encoder from the ROB head, and is
marked invalid when no branch is pending. A branch that resolves in one cycle
releases its dependent transmitters in the next.

Every cycle, every table entry whose YRoT has become safe is cleared. This eager 
clearing keeps every tag valid as a ROB index. A root can only retire from the 
ROB head, where every unresolved branch is younger than it, so its tag is always 
cleared before its ROB slot can be reused. 
It also means a branch mispredict needs no snapshot of the taint tables. 
Surviving registers keep their taint and continue to clear, and a register freed 
by a squashed instruction returns to BOOM's free list, where dispatch overwrites 
its taint when the register is next allocated. A full flush clears both tables, since
committed values are untainted. During recovery, \texttt{stt\_block} holds
every transmitter. It is set for either phase of a mispredict, during a
rollback, and for the two cycles after a flush.

% Issue #5 (A3): Age Matrix.
\paragraph{The $\mu$STT Age Matrix}
In baseline STT, the same-cycle bypass chains from lane to lane so each lane's
YRoT can depend on the result of every older lane in the group, and its delay
grows with dispatch width. The \emph{Age Matrix} breaks this chain. It
compares the taint read from the table for every source in the group pairwise,
in parallel with the dependency check, so no comparison waits for an older
lane's result. Each lane then finds the older lanes it reaches through
same-cycle dependencies and selects the youngest table taint among their
sources. If a lane reaches a load in its own group, that load's own ROB index
wins without a comparison, since a same-group root is younger than any taint
already in the table. We port $\mu$STT's Verilog to Chisel at BOOM's widths,
behind \texttt{enableAgeMatrix}, sharing the tables, clearing and writes with
the bypass. A structural change simplifies the port from $\mu$STT. 
Instead of computing two separate dependency closures, ours computes a single one 
(reachability through lanes that are not loads) and selects exactly the same set of
taints. Building the bypass also exposed a scaling problem. Verilog
generation inlined each lane's taint into every younger lane, so the
expression grew exponentially with width, and at width 4 Verilator could not
compile it. Keeping each lane's taint as a named wire fixes this without
changing the logic.

% Owner: @SamirJoshiGit (track B). Only the owner edits this file; send input to them on #40. Transmitter delay, LSU audit and the impede micro-op.

% Issue #6 (B1): conservative stand-in; #7 (B2): real STT delay + LSU audit.
\paragraph{Delaying transmitters}
Without waiting for taint tracking, we first build a \emph{conservative stand-in}
which blocks every transmitter from issuing while any older branch is unresolved.
BOOM already tracks these branches in the branch mask of each micro-op, so this needs no new 
state and enables the policy which is strictly stronger than STT. Crucially, the policy lets 
security testing begin early and gives an upper bound on the performance cost. After building the 
conservative stand-in, we replace the branch-mask condition with the \textit{isSafe(YRoT)} 
check; however, blocking issue at the issue queue is not enough on its own. BOOM's load/\allowbreak 
store unit can fire loads without a fresh issue because of load-queue retries, speculative wakeup, 
TLB-miss replay, store-to-load forwarding, and the prefetcher, which can all bypass the issue
queue. As a result, we must audit each of these paths so that all of them respect taint.

% Issue #8 (B3): impede micro-op.
\paragraph{The $\mu$STT impede micro-op}
The \emph{impede micro-op} is a special 
micro-op that removes the per-entry delay logic from the issue queue by sitting in front of each transmitter 
and waiting until the VP passes its YRoT. The transmitter depends on the special 
micro-op, so BOOM's ordinary wakeup logic will delay it. In gem5, the ARM model inserts 
this micro-op by splitting instructions at decode, which BOOM does not do. We therefore evaluate 
either emitting an extra micro-op at dispatch, or reserving a dedicated issue-queue entry to hold 
the special micro-op.

\subsection{Evaluation}

% Owner: @vinle0 (track C). Only the owner edits this file; send input to them on #40. How security and performance are measured; the numbers go in 6a.

% Issues #9 (C1), #10 (C2), #11 (C3).
\paragraph{Security and performance (track C)}
Security is tested with the BOOM Spectre proofs of concept, ported to Chipyard 
bare-metal tests, together with the functional and security checks of 
Section~\ref{sec:verification}. Performance is measured as cycles and IPC from RTL simulation in Verilator, 
using CoreMark, Dhrystone, and the riscv-tests benchmarks, and normalized to the baseline.

% Owner: @erickan2002-del (track D). Only the owner edits this file; send input to them on #40. How hardware cost is measured; the numbers go in 6b.

% Issues #12 (D1), #13 (D2), #14 (D3), #15 (D4), #16 (D5).
\paragraph{Hardware cost (track D)}
We have no FPGA board available currently, so hardware cost comes from Vivado 2023.2 targeting the
Xilinx VCU118 (\texttt{xcvu9p}), using Chipyard's existing BOOM FPGA configuration. The full flow 
runs synthesis through implementation. It gives hierarchical utilization (to attribute logic to rename, 
issue, ROB and LSU), timing slack, and power estimates. We utilize out-of-context synthesis of the BOOM 
tile alone to provide a faster loop for checking design changes. We also synthesize $\mu$STT's released 
rename-stage RTL at BOOM's parameters to reproduce the paper's area and delay claims and give a reference 
for our Chisel versions.


\subsection{Integration}
Table~\ref{tab:configs} lists the configurations that combine these pieces.
Each is a Chipyard configuration in two sizes, Small (one instruction
dispatched per cycle) and Medium (two), named
\texttt{SmallBoomV3}\textit{X}\texttt{Config} and
\texttt{MediumBoomV3}\textit{X}\texttt{Config}; Medium is the primary
configuration for development and results. Each is measured on all three axes
(security, performance and hardware cost) against the unmodified Chipyard~1.15.0
with SonicBOOM, which we froze as the baseline.

\begin{table}[t]
  \caption{Configurations evaluated.}
  \label{tab:configs}
  \small
  \begin{tabular}{@{}lll@{}}
    \toprule
    Configuration & Taint tracking & Delay (issue/LSU) \\
    \midrule
    Baseline          & none             & none \\
    Conservative      & none             & any unresolved branch \\
    STT               & bypass\textsuperscript{\dag}     & \textit{isSafe(YRoT)} \\
    STT + Age Matrix  & Age Matrix\textsuperscript{\dag} & \textit{isSafe(YRoT)} \\
    $\mu$STT          & Age Matrix\textsuperscript{\dag} & impede micro-op \\
    \bottomrule
  \end{tabular}

  \smallskip
  {\footnotesize \textsuperscript{\dag}Includes the taint table.}
\end{table}
